\NativeSlot{43}{77}{%
% p103_s01

\ExamQuestion{6（续）}{5}%
\par\smallskip\Needspace{4\baselineskip}\begingroup\small\setlength{\tabcolsep}{3pt}
\noindent\begin{tabularx}{\linewidth}{|*{4}{>{\centering\arraybackslash}X|}}\hline
 & 第一名 & 第二名 & 第三名\\\hline
一组 & A：15秒 & B：20秒 & C：23秒\\\hline
二组 & 甲：16秒 & 乙：21秒 & 丙：24秒\\\hline
\end{tabularx}\par
\endgroup\smallskip
\ExamOptions{A. 丙—A，乙—B，甲—C}{B. 丙—B，乙—C，甲—A}{C. 丙—A，乙—C，甲—B}{D. 丙—C，乙—A，甲—B}
}
\NativeSlot{178}{76}{%
% p103_s02

\ExamQuestion{7.}{5}%
小军在探索四边形的内角和时，采用了如下图所示的分割转化法，下列选项中，\mbox{（\hspace{1.8em}）}体现了他的推算过程。\par
\ExamOptions{A. \(4\times180^\circ-360^\circ\)}{B. \(\angle1+\angle2+\angle3+\angle4\)}{C. \(3\times180^\circ-180^\circ\)}{}
\BookFigure{figures/p103_s02-01.tex}
}
\NativeSlot{306}{54}{%
\ExamGroup{三、计算。{\normalfont\kaishu\fontsize{10.8bp}{14bp}\selectfont （共30分）}}

% p103_s03

\ExamQuestion{1.}{5}%
直接写得数。（8分）\par
\noindent \NativeOralCell{3.26+0.4=}\NativeOralCell{4.37+2.7=}\NativeOralCell{5.6-4=}\NativeOralCell{62+38\times0=}\par
\noindent \NativeOralCell{0.65+0.35=}\NativeOralCell{0.6\div100=}\NativeOralCell{8-2.68=}\NativeOralCell{4.32\times1000=}\par
}
\NativeSlot{404}{143}{%
\par\noindent 2.\quad 列竖式计算。（带★的要验算）（13分）\par
\par\noindent%
\ExamTile{.242}{%
% p103_s04

\(9.27+15.9=\hspace{1.8em}\)\par

\ExamWork{7}
}\hfill%
\ExamTile{.242}{%
% p103_s05

\(25.7+16.4=\hspace{1.8em}\)\par

\ExamWork{5}
}\hfill%
\ExamTile{.242}{%
% p103_s06

\(10-4.03=\hspace{1.8em}\)\par

\ExamWork{5}
}\hfill%
\ExamTile{.242}{%
% p103_s07

★\(18.2-6.78=\hspace{1.8em}\)\par

\ExamWork{7}
}\par

\par\noindent 3.\quad 计算下列各题，能简便计算的要简便计算。（9分）\par
\par\noindent%
\ExamTile{.323}{%
% p103_s08

\(45.01+6.1+0.99+2.9\)\par

\ExamWork{3}
}\hfill%
\ExamTile{.323}{%
% p103_s09

\(45\times98\)\par

\ExamWork{3}
}\hfill%
\ExamTile{.323}{%
% p103_s10

\(101\times101-101\)\par

\ExamWork{3}
}\par
}
\NativeSlot{641}{90}{%
\ExamGroup{四、操作。{\normalfont\kaishu\fontsize{10.8bp}{14bp}\selectfont （共12分）}}

\par\noindent%
\ExamTile{.54}{%
% p103_s11

\ExamQuestion{（1）}{5}%
画出三角形\(ABC\)中\(AC\)边上的高。\par

% p103_s12

\ExamQuestion{（2）}{5}%
以\(AB\)边所在直线为对称轴，画出轴对称图形的另一半。\par

% p103_s13

\ExamQuestion{（3）}{5}%
把三角形\(ABC\)先绕点\(A\)逆时针旋转\(90^\circ\)，再向右平移\(4\)格。\par
}\hfill%
\ExamTile{.435}{%
\BookFigure{figures/p103_s11-01.tex}
}\par
}
